\textbf{COLUNI 2025 -- Questão 24}\par\medskip Três estudantes, Tiago, Flávio e Josy, realizaram, no laboratório de ciências, o experimento de eletrização, apresentado abaixo, em que um canudinho de plástico é atritado em uma toalha de papel.

\begin{center}
\setlength{\unitlength}{1mm}
\begin{picture}(100,52)
\put(8,18){\line(3,1){42}}
\put(20,30){$-$}\put(28,33){$-$}\put(36,36){$-$}\put(44,39){$-$}\put(52,42){$-$}\put(58,39){$-$}
\put(64,10){\framebox(25,36){}}
\put(68,37){$+$}\put(76,37){$+$}\put(84,37){$+$}\put(68,27){$+$}\put(76,27){$+$}\put(84,27){$+$}
\end{picture}
\end{center}

Considerando que, após o processo, o canudinho encontra-se eletrizado negativamente e a toalha de papel encontra-se eletrizada positivamente, os estudantes fizeram anotações sobre o experimento com as seguintes afirmações:

\begin{enumerate}
\item[Josy:] O canudinho recebeu elétrons dos átomos existentes na toalha de papel.
\item[Flávio:] A toalha de papel recebeu prótons dos átomos existentes no canudinho.
\item[Tiago:] A toalha de papel transferiu elétrons de seus átomos para o canudinho.
\end{enumerate}
Está(ão) CORRETA(S) a(s) afirmativa(s) de:
\begin{enumerate}
\item Josy e Flávio.
\item Flávio e Tiago.
\item Josy e Tiago.
\item Josy, apenas.
\end{enumerate}
